% Lecture 01 — Limits, an example didactic document.
%
% Build:  lualatex lecture-01-limits.tex
% Picks up coursemeta.tex from this directory automatically.

\documentclass[
	unit   = Lecture,
	number = 1,
	title  = {Limits and Continuity},
]{didactic}

% Overview is prose, so it can't ride in the class option list — set it with the
% \overview content-setter instead.
\overview{%
	Intuitive definition of the limit, one-sided limits, and a
	first look at continuity via the Intermediate Value Theorem.
}

\begin{document}

\maketitle

\section{Limits}

\begin{definition}[Limit]
	We write $\lim_{x \to a} f(x) = L$ to mean that $f(x)$ can be
	made arbitrarily close to $L$ by taking $x$ sufficiently close
	to (but not equal to) $a$.
\end{definition}

\begin{example}
	$\displaystyle \lim_{x \to 3} (2x + 1) = 7$.
\end{example}

\section{Continuity}

\begin{definition}[Continuity at a point]
	A function $f$ is continuous at $a$ if
	$\lim_{x \to a} f(x) = f(a)$.
\end{definition}

\begin{theorem}[Intermediate Value]
	If $f$ is continuous on $[a, b]$ and $N$ lies between $f(a)$
	and $f(b)$, then there exists $c \in (a, b)$ with $f(c) = N$.
\end{theorem}

The \ivt{} is the workhorse behind the following exercise. (\texttt{\textbackslash ivt}
is defined once in \texttt{course-preamble.tex}, auto-loaded for every document in
this course via coursemeta's \texttt{preamble-file} key.)

\begin{exercise}
	Show that $x^3 - x - 1 = 0$ has a solution in $(1, 2)$.
\end{exercise}

\hint Evaluate $f$ at the endpoints $x = 1$ and $x = 2$ and compare the signs.

\begin{solution}[2in]
	Let $f(x) = x^3 - x - 1$. Then $f(1) = -1 < 0$ and
	$f(2) = 5 > 0$. Since $f$ is continuous, the IVT gives a root
	in $(1, 2)$.
\end{solution}

\begin{readings}
	\item Stewart, \emph{Calculus}, \S2.5 (Continuity).
	\item Course notes, Lecture 1.
\end{readings}

\end{document}
